% problem_id: S001-lemma-pre-derived-chern-class
% source_id: S001 (Stacks Project)
% source_file: chow.tex
% statement_locator: chow.tex lines 9090--9104
% proof_locator: chow.tex lines 9106--9210
% env: lemma
% pairing: tight (verified)
% status: complete (statement + author proof verbatim + reason + locators)
%
\begin{lemma}
\label{lemma-pre-derived-chern-class}
Let $(S, \delta)$ be as in Situation \ref{situation-setup}.
Let $X$ be locally of finite type over $S$. Let $E \in D(\mathcal{O}_X)$
be an object such that there exists a locally bounded complex
$\mathcal{E}^\bullet$ of finite locally free $\mathcal{O}_X$-modules
representing $E$. Then a slight generalization of the above constructions
$$
c(\mathcal{E}^\bullet) \in \prod\nolimits_{p \geq 0} A^p(X),\quad
ch(\mathcal{E}^\bullet) \in
\prod\nolimits_{p \geq 0} A^p(X) \otimes \mathbf{Q},\quad
P_p(\mathcal{E}^\bullet) \in A^p(X)
$$
are independent of the choice of the complex $\mathcal{E}^\bullet$.
\end{lemma}

\begin{proof}
We prove this for the total Chern class; the other two cases follow
by the same arguments using
Lemma \ref{lemma-chern-character-additive}
instead of
Lemma \ref{lemma-additivity-chern-classes}.

\medskip\noindent
As in Remark \ref{remark-extend-to-finite-locally-free} in order
to define the total chern class $c(\mathcal{E}^\bullet)$
we decompose $X$ into open and closed subschemes
$$
X = \coprod\nolimits_{i \in I} X_i
$$
such that the rank $\mathcal{E}^n$ is constant on $X_i$ for
all $n$ and $i$. (Since these ranks are locally constant functions
on $X$ we can do this.) Since $\mathcal{E}^\bullet$ is locally
bounded, we see that only a finite number of the sheaves
$\mathcal{E}^n|_{X_i}$ are nonzero for a fixed $i$. Hence we
can define
$$
c(\mathcal{E}^\bullet|_{X_i}) =
\prod\nolimits_n c(\mathcal{E}^n|_{X_i})^{(-1)^n}
\in \prod\nolimits_{p \geq 0} A^p(X_i)
$$
as above. By Lemma \ref{lemma-disjoint-decomposition-bivariant}
we have $A^p(X) = \prod_i A^p(X_i)$. Hence for each $p \in \mathbf{Z}$
we have a unique element $c_p(\mathcal{E}^\bullet) \in A^p(X)$ restricting
to $c_p(\mathcal{E}^\bullet|_{X_i})$ on $X_i$ for all $i$.

\medskip\noindent
Suppose we have a second locally bounded complex
$\mathcal{F}^\bullet$ of finite locally free $\mathcal{O}_X$-modules
representing $E$.
Let $g : Y \to X$ be a morphism locally of finite type with $Y$ integral.
By Lemma \ref{lemma-bivariant-zero} it suffices to show that
with $c(g^*\mathcal{E}^\bullet) \cap [Y]$ is the same as
$c(g^*\mathcal{F}^\bullet) \cap [Y]$ and it even suffices to prove
this after replacing $Y$ by an integral scheme proper and birational
over $Y$. Then first we conclude that $g^*\mathcal{E}^\bullet$
and $g^*\mathcal{F}^\bullet$ are bounded complexes of finite locally
free $\mathcal{O}_Y$-modules of constant rank. Next, by
More on Flatness, Lemma \ref{flat-lemma-blowup-complex-integral}
we may assume that $H^i(Lg^*E)$ is perfect of tor dimension $\leq 1$
for all $i \in \mathbf{Z}$.
This reduces us to the case discussed in the next paragraph.

\medskip\noindent
Assume $X$ is integral, $\mathcal{E}^\bullet$ and $\mathcal{F}^\bullet$
are bounded complexes of finite locally free modules of constant rank, and
$H^i(E)$ is a perfect $\mathcal{O}_X$-module of tor dimension $\leq 1$
for all $i \in \mathbf{Z}$. We have to
show that $c(\mathcal{E}^\bullet) \cap [X]$ is the same as
$c(\mathcal{F}^\bullet) \cap [X]$. Denote
$d_\mathcal{E}^i : \mathcal{E}^i \to \mathcal{E}^{i + 1}$ and
$d_\mathcal{F}^i : \mathcal{F}^i \to \mathcal{F}^{i + 1}$
the differentials of our complexes. By
More on Flatness, Remark \ref{flat-remark-when-you-have-a-complex}
we know that $\Im(d_\mathcal{E}^i)$, $\Ker(d_\mathcal{E}^i)$,
$\Im(d_\mathcal{F}^i)$, and $\Ker(d_\mathcal{F}^i)$
are finite locally free $\mathcal{O}_X$-modules for all $i$.
By additivity (Lemma \ref{lemma-additivity-chern-classes}) we see that
$$
c(\mathcal{E}^\bullet) = \prod\nolimits_i
c(\Ker(d_\mathcal{E}^i))^{(-1)^i} c(\Im(d_\mathcal{E}^i))^{(-1)^i}
$$
and similarly for $\mathcal{F}^\bullet$. Since we have the
short exact sequences
$$
0 \to \Im(d_\mathcal{E}^i) \to \Ker(d_\mathcal{E}^i) \to H^i(E) \to 0
\quad\text{and}\quad
0 \to \Im(d_\mathcal{F}^i) \to \Ker(d_\mathcal{F}^i) \to H^i(E) \to 0
$$
we reduce to the problem stated and solved in the next paragraph.

\medskip\noindent
Assume $X$ is integral and we have two short exact sequences
$$
0 \to \mathcal{E}' \to \mathcal{E} \to \mathcal{Q} \to 0
\quad\text{and}\quad
0 \to \mathcal{F}' \to \mathcal{F} \to \mathcal{Q} \to 0
$$
with $\mathcal{E}$, $\mathcal{E}'$, $\mathcal{F}$, $\mathcal{F}'$
finite locally free. Problem: show that
$c(\mathcal{E})c(\mathcal{E}')^{-1} \cap [X] =
c(\mathcal{F})c(\mathcal{F}')^{-1} \cap [X]$.
To do this, consider the short exact sequence
$$
0 \to \mathcal{G} \to \mathcal{E} \oplus \mathcal{F} \to \mathcal{Q} \to 0
$$
defining $\mathcal{G}$. Since $\mathcal{Q}$ has tor dimension $\leq 1$
we see that $\mathcal{G}$ is finite locally free. A diagram chase
shows that the kernel of the surjection $\mathcal{G} \to \mathcal{F}$
maps isomorphically to $\mathcal{E}'$ in $\mathcal{E}$ and
the kernel of the surjection $\mathcal{G} \to \mathcal{E}$ maps
isomorphically to $\mathcal{F}'$ in $\mathcal{F}$. (Working affine
locally this follows from or is equivalent to Schanuel's lemma, see
Algebra, Lemma \ref{algebra-lemma-Schanuel}.)
We conclude that
$$
c(\mathcal{E})c(\mathcal{F}') = c(\mathcal{G}) =
c(\mathcal{F})c(\mathcal{E}')
$$
as desired.
\end{proof}
